The fourth question was whether the loop reaches sound, auditable conclusions on a real interpretability question when the executor runs its own code. We asked whether Geneformer attention encodes transcription-factor (TF)$\to$target regulation beyond co-expression, expression-rank proximity and hub (degree) structure, using a data kit built from 1,000 Tabula Sapiens immune cells (Materials and methods; S3 File). Whether the agents answered this question correctly can be judged only against a known answer, so we first report our own analysis of the kit and then the agents' executed runs.

\subsubsection*{Independent analysis}

Our own analysis of the kit serves as the reference against which the agents' executed analyses are judged (Fig~\ref{fig5}; S3 File).

\textbf{Geneformer attention is not symmetric.} Across 873,181 gene pairs co-present in at least 50 cells, the head-averaged attention from gene $i$ to gene $j$ was only weakly correlated with the attention from $j$ to $i$ in every layer (Spearman $\rho$ 0.16--0.32; median relative asymmetry $|A_{ij}-A_{ji}|/(A_{ij}+A_{ji})$ 0.14--0.31), and for 11--46\% of pairs one direction exceeded the other at least twofold (Fig~\ref{fig5}A). Individual heads were less symmetric still (median $\rho$ 0.10 across the 144 heads, range --0.22 to 0.75). This is what the architecture implies: queries and keys use different learned projections, and each query's attention is normalised by its own row-wise softmax~\cite{vaswani2017attention,tsai2019dissection}, so $A_{ij}=A_{ji}$ is neither guaranteed nor observed. On TRRUST edges the TF-as-query direction received more attention than the target-as-query direction in several later layers (layer 12 median $\log_2$ ratio 0.53, 95\% CI 0.35--0.75), but TF--gene pairs that are not edges showed a similar bias in most layers, so the direction of attention mostly reflects which token is a TF, not which pairs are regulatory.

\textbf{Apparent edge recovery is mostly hub structure.} Among 310,435 ordered TF--gene candidate pairs (235 TRRUST TFs; 424 TRRUST edges, base rate 0.14\%), symmetrised attention ranked edges above non-edges in every layer, best in layer 3 (AUROC 0.673, 95\% CI 0.631--0.708; AUPRC 0.0037, 2.7 times the base rate), and above absolute co-expression (AUROC 0.545, 0.511--0.575) in layers 1--7 (Benjamini--Hochberg $q<0.05$). Rank proximity explained little of attention (Spearman $\rho$ with rank distance --0.29 to --0.05 across layers). The decisive control was the degree-preserving null, which rewires the reference edges while keeping every TF's number of targets and every target's number of regulators: it alone reached an AUROC of 0.645 in layer 3 and reproduced 64--102\% of attention's AUROC above 0.5 across layers. Attention still exceeded this null by a small margin in six of the 12 layers (0.017--0.046 AUROC; Benjamini--Hochberg $q<0.05$; layer 3: 0.673 vs 0.645, max-statistic $p=0.004$ over layers). A target's in-degree in TRRUST, computed without the scored edge, reached an AUROC of 0.829 by itself. In post hoc decompositions, the gene-level component of attention (how much a gene attends and is attended to overall) matched the null exactly (layer 3: 0.654 vs 0.653), and attention's advantage over co-expression was no larger than the advantage expected from degree structure in any layer (0 of 36 layer-by-orientation tests). A pair-specific component remained against the degree null in layers 1--6 (layer 3: AUROC 0.592 vs null 0.551, $q=0.006$), but after further adjustment for co-expression, expression and rank distance it was not detectable with TRRUST (0 of 36 tests; best 0.612 vs 0.575, $q=0.07$) and was detectable with the larger DoRothEA A--C reference (24 of 36 tests), with an excess of at most about 0.04 AUROC over the null in either reference.

\textbf{Interpretation.} The direct evidence, from the pre-specified tests, is that reference degree structure alone reproduces most of the agreement between Geneformer attention and curated regulatory edges, although attention retains a small, significant excess over the degree-preserving null. Our inference, which rests on the post hoc decompositions, is that this agreement is driven mainly by which genes are hubs in the reference and which genes attend or are attended to strongly, rather than by the specific TF--target pairs, and that a small pair-specific component exists but is close to the detection limit of the available references; the difference between TRRUST and DoRothEA is consistent with the larger reference giving more power, although DoRothEA's lower-confidence levels derive partly from expression-based inference, so residual non-linear co-expression could contribute. A plausible but untested explanation for the hub-level agreement is that broadly expressed, well-studied genes are both over-annotated in curated databases and preferentially attended. For the purpose of this study, these results define what a correct automated analysis of the kit should conclude: raw edge recovery overstates regulatory information, and any claim of TF--target encoding must survive a degree-preserving null. Although symmetry was not part of the question, no correct analysis can treat attention as symmetric.
